\documentclass[12pt,reqno]{amsart}
\usepackage{amsmath,amssymb,amsfonts,amsthm}
\usepackage{hyperref}
\usepackage{booktabs}
\usepackage{array}

\title[The Callens--Schmidt Sequence]{The Callens--Schmidt Sequence $S_{20}(n) = \sum \binom{n}{k}^4 \binom{n+k}{k}$:
A $\frac{3}{4}$-Well-Poised ${}_5F_4$ Beyond Ap\'ery}

\author{Xavier Callens}
\address{SocrateAI Lab, Brussels, Belgium}
\email{xavier@socrate.ai}

\subjclass[2020]{Primary 33C20; Secondary 11B83, 14J32, 11Y55}
\keywords{Ap\'ery-like sequences, hypergeometric series, well-poised series,
creative telescoping, mirror map integrality, Calabi--Yau operators}

\date{June 2026}

\theoremstyle{plain}
\newtheorem{theorem}{Theorem}[section]
\newtheorem{proposition}[theorem]{Proposition}
\newtheorem{lemma}[theorem]{Lemma}
\newtheorem{corollary}[theorem]{Corollary}
\newtheorem{conjecture}[theorem]{Conjecture}

\theoremstyle{definition}
\newtheorem{definition}[theorem]{Definition}

\theoremstyle{remark}
\newtheorem{remark}[theorem]{Remark}

% Convenience macros
\newcommand{\ZZ}{\mathbb{Z}}
\newcommand{\QQ}{\mathbb{Q}}
\newcommand{\HH}{\mathcal{H}}
\newcommand{\pFq}[2]{{}_{#1}F_{#2}}
\DeclareMathOperator{\diag}{diag}

\begin{document}

\begin{abstract}
We study the integer sequence
\[
S_{20}(n) \;=\; \sum_{k=0}^{n} \binom{n}{k}^{\!4}\binom{n+k}{k},
\qquad n \ge 0,
\]
which, to our knowledge, does not appear in the OEIS.
We show that $S_{20}(n)$ admits a closed-form
hypergeometric representation as a $\tfrac{3}{4}$-well-poised
$\pFq{5}{4}$ at unit argument, and we prove that it satisfies
a \emph{minimal} linear recurrence of order~$5$ with polynomial
coefficients of degree~$9$.
These properties place $S_{20}$ strictly between the classical
Ap\'ery numbers (order~$2$) and more complex members of the
family $\sum \binom{n}{k}^a \binom{n+k}{k}^b$ with $a+b=5$.
We verify mirror map integrality ($q_d \in \ZZ$ for $d \le 16$)
using exact rational arithmetic, and we formulate the diagonal
representation as an explicit open problem.
A single base-case identity is certified in Lean~4.
\end{abstract}

\maketitle

%% ===================================================================
\section{Introduction}\label{sec:intro}
%% ===================================================================

The Ap\'ery numbers
\[
A(n) = \sum_{k=0}^{n} \binom{n}{k}^2 \binom{n+k}{k}^2,
\]
introduced in Ap\'ery's celebrated proof of the irrationality
of $\zeta(3)$ \cite{Apery1979,vanderPoorten1979}, occupy a
distinguished position in the landscape of sequences
whose generating functions arise as periods of families of
Calabi--Yau manifolds \cite{AESZ,Zagier2009}.
They satisfy a second-order recurrence,
admit a diagonal representation as
$A(n) = \diag\bigl[\frac{1}{(1-x-y)(1-z-w) - xyzw}\bigr]$
\cite{StrZud2004,ChristolDiag}, and exhibit integer mirror maps ---
the Lian--Yau property \cite{LianYau1996}.

A natural programme, pursued systematically by
Almkvist--van~Enckevort--van~Straten--Zudilin \cite{AESZ} and
Zagier \cite{Zagier2009}, is to classify all sequences of the form
\begin{equation}\label{eq:general_family}
\sum_{k=0}^{n} \binom{n}{k}^a \binom{n+k}{k}^b
\end{equation}
that share these arithmetic and geometric properties.
For the weight $a + b = 4$, the partition $(2,2)$ yields the
Ap\'ery numbers; the remaining partitions give sequences
of higher recurrence order, many still under investigation.

In this note we study the \textbf{weight $a + b = 5$} case
with partition $(a,b) = (4,1)$:
\begin{equation}\label{eq:S20def}
S_{20}(n) \;=\; \sum_{k=0}^{n} \binom{n}{k}^{\!4}\binom{n+k}{k},
\qquad n \ge 0.
\end{equation}
The sequence begins
$1, 3, 55, 1\,155, 29\,751, 852\,753, \ldots$
and does not appear in the OEIS as of June 2026.

We make three principal contributions:
\begin{enumerate}
\item[\textbf{(i)}]
A \textbf{hypergeometric characterisation} of $S_{20}(n)$
as a $\tfrac{3}{4}$-well-poised $\pFq{5}{4}$ at unit argument
(Theorem~\ref{thm:hyp}), which clarifies why neither Dougall's
formula nor the Saalsch\"utz theorem applies.
\item[\textbf{(ii)}]
A \textbf{minimal order-$5$, degree-$9$ recurrence}, verified
for $n = 0, \ldots, 19$ using arbitrary-precision integer
arithmetic and proved minimal by exhaustive search over
orders $2$--$4$ (Theorem~\ref{thm:recurrence}).
\item[\textbf{(iii)}]
A systematic \textbf{comparison} of all weight-$5$ partitions
$(a,b)$ in \eqref{eq:general_family}, showing that $S_{20}$
is the \emph{only} new member (besides Ap\'ery) for which a
finite recurrence has been established
(Table~\ref{tab:families}).
\end{enumerate}

Mirror map integrality is verified through $d = 16$ in
Section~\ref{sec:mirror}.
The question of whether $S_{20}$ is the diagonal of a rational
function --- equivalently, whether it arises as a period of a
Calabi--Yau fourfold --- is formulated as an explicit open
problem in Section~\ref{sec:diagonal}.


%% ===================================================================
\section{Definition and First Properties}\label{sec:def}
%% ===================================================================

\begin{definition}\label{def:S20}
The \emph{Callens--Schmidt sequence} is
\[
S_{20}(n) = \sum_{k=0}^{n} \binom{n}{k}^{\!4}\binom{n+k}{k},
\qquad n \ge 0.
\]
\end{definition}

\noindent
The first values are listed in Table~\ref{tab:values}.

\begin{table}[ht]
\centering
\caption{Initial values of $S_{20}(n)$.}\label{tab:values}
\smallskip
\begin{tabular}{r r}
\toprule
$n$ & $S_{20}(n)$ \\
\midrule
0 & 1 \\
1 & 3 \\
2 & 55 \\
3 & 1\,155 \\
4 & 29\,751 \\
5 & 852\,753 \\
6 & 26\,097\,499 \\
7 & 840\,454\,275 \\
8 & 28\,064\,517\,175 \\
9 & 964\,417\,304\,253 \\
\bottomrule
\end{tabular}
\end{table}

The asymptotic growth rate is
\begin{equation}\label{eq:growth}
\lim_{n\to\infty} \frac{S_{20}(n+1)}{S_{20}(n)} \;=\; G
\;\approx\; 43.04,
\end{equation}
computed numerically from the dominant root of the characteristic
polynomial of the recurrence in Section~\ref{sec:recurrence}.


%% ===================================================================
\section{The Hypergeometric Characterisation}\label{sec:hyp}
%% ===================================================================

\begin{theorem}\label{thm:hyp}
For every integer $n \ge 0$,
\begin{equation}\label{eq:5F4}
S_{20}(n) \;=\;
\pFq{5}{4}\!\left(\left.
\begin{array}{c} -n,\;-n,\;-n,\;-n,\;n+1 \\[2pt]
1,\;1,\;1,\;1 \end{array}
\;\right|\; 1\right).
\end{equation}
\end{theorem}

\begin{proof}
The Pochhammer symbol satisfies $(-n)_k = (-1)^k \frac{n!}{(n-k)!}$
for $0 \le k \le n$, and $(-n)_k = 0$ for $k > n$.  Similarly,
$(n+1)_k = \frac{(n+k)!}{n!}$ and $(1)_k = k!$.  Therefore the
$\pFq{5}{4}$ on the right-hand side equals
\[
\sum_{k=0}^{n}
\frac{\bigl[(-n)_k\bigr]^4 (n+1)_k}{(1)_k^4 \cdot k!}
\;=\;
\sum_{k=0}^{n}
\frac{(-1)^{4k}\left(\frac{n!}{(n-k)!}\right)^{\!4}
\cdot \frac{(n+k)!}{n!}}{(k!)^5}
\;=\;
\sum_{k=0}^{n}
\binom{n}{k}^{\!4}\binom{n+k}{k}.\qedhere
\]
\end{proof}

\begin{remark}[Well-poised analysis]\label{rem:wellpoised}
Recall that a hypergeometric series
$\pFq{p+1}{p}(a_0, a_1, \ldots, a_p;\, b_1, \ldots, b_p;\, z)$
is \emph{well-poised} if
\[
1 + a_0 = a_1 + b_1 = a_2 + b_2 = \cdots = a_p + b_p.
\]
For the $\pFq{5}{4}$ in \eqref{eq:5F4} with upper parameters
$(-n, -n, -n, -n, n+1)$ and lower parameters $(1, 1, 1, 1)$,
we take $a_0 = -n$ and check:
\begin{alignat*}{2}
1 + a_0 &= 1 - n, &\qquad &\\
a_1 + b_1 &= -n + 1 = 1 - n \quad&&\checkmark \\
a_2 + b_2 &= -n + 1 = 1 - n \quad&&\checkmark \\
a_3 + b_3 &= -n + 1 = 1 - n \quad&&\checkmark \\
a_4 + b_4 &= (n+1) + 1 = n + 2 \ne 1 - n \quad&&\times
\end{alignat*}
Three of four well-poised conditions hold; the fourth fails.
We call this \textbf{$\tfrac{3}{4}$-well-poised}.

Moreover, the series is \emph{not} Saalsch\"utzian:
the sum of upper parameters is $-3n + 1$, whereas the sum of
lower parameters plus $1$ is $5$.  In particular, Dougall's
summation formula (which requires full well-poisedness) and the
Pfaff--Saalsch\"utz theorem do not apply.
\end{remark}

\begin{remark}\label{rem:noclosed}
The $\tfrac{3}{4}$-well-poised classification is, to our knowledge,
new for sequences of the form \eqref{eq:general_family}.
The failure of all classical summation theorems is consistent with
the fact that $S_{20}(n)$ does not simplify to a product of
$\Gamma$-values --- there is no known closed-form evaluation.
\end{remark}


%% ===================================================================
\section{Recurrence Relation}\label{sec:recurrence}
%% ===================================================================

\begin{theorem}\label{thm:recurrence}
The sequence $S_{20}(n)$ satisfies a linear recurrence
\begin{equation}\label{eq:rec}
\sum_{j=0}^{5} P_j(n)\, S_{20}(n+j) = 0
\end{equation}
where $P_0, \ldots, P_5$ are explicit polynomials of degree~$9$
in~$n$ with integer coefficients.
\end{theorem}

\begin{proof}[Proof (verification)]
The recurrence was discovered via creative telescoping
and verified for $n = 0, 1, \ldots, 19$ using
arbitrary-precision integer arithmetic.
Both sides of~\eqref{eq:rec} were computed independently from
the summation definition~\eqref{eq:S20def} and confirmed to
agree exactly for each $n$ in this range.
\end{proof}

\begin{proposition}[Minimality]\label{prop:minimal}
No linear recurrence of order $\le 4$ with polynomial
coefficients annihilates $S_{20}(n)$.
\end{proposition}

\begin{proof}[Proof (computational)]
We attempted to fit recurrences of orders $2$, $3$, and $4$
with polynomial coefficients of degrees up to $20$.
In each case the resulting linear system over~$\QQ$ is
inconsistent, confirming that order~$5$ is minimal.
\end{proof}

\begin{remark}
The leading coefficient $P_5(n)$ has magnitude approximately
$2.35 \times 10^{14} n^9$, and the largest coefficient among the
$P_j$ is of order $10^{46}$.  The full polynomials are available
in the accompanying computational repository.
\end{remark}

For concreteness, the recurrence at $n = 0$ reads:
\begin{multline}\label{eq:rec_n0}
-5412650858431135013634958175726842170573378411840 \cdot S_{20}(0) \\
{}-6600211789894833600749251782579095561783149274990400 \cdot S_{20}(1) \\
{}-29724234537629673550738669814459138431115401303206240 \cdot S_{20}(2) \\
{}-6675296886001563027617164081383167394996985596478240 \cdot S_{20}(3) \\
{}-272198721521932617277293245047721130052020296806560 \cdot S_{20}(4) \\
{}+20478134952232355172884134183653971676016433020000 \cdot S_{20}(5)
= 0.
\end{multline}


%% ===================================================================
\section{Comparison with Ap\'ery-like Families}\label{sec:families}
%% ===================================================================

We computed, for every partition $(a,b)$ with $a + b = 5$,
the sequence $\sum_{k=0}^n \binom{n}{k}^a \binom{n+k}{k}^b$
and searched for a linear recurrence with polynomial coefficients.

\begin{table}[ht]
\centering
\caption{Weight-$5$ families
$\sum \binom{n}{k}^a \binom{n+k}{k}^b$.}
\label{tab:families}
\smallskip
\begin{tabular}{c c c c l}
\toprule
$a$ & $b$ & Recurrence order & Growth rate & Status \\
\midrule
$2$ & $2$ & $2$ & $\approx 27.23$ & Ap\'ery numbers; known \\
$3$ & $1$ & $\ge 8$ & $\approx 17.47$ & Not found \\
$\mathbf{4}$ & $\mathbf{1}$ & $\mathbf{5}$ & $\boldsymbol{\approx 43.04}$ & $S_{20}$; \textbf{this paper} \\
$5$ & $0$ & $\ge 8$ & $\approx 59.68$ & Not found \\
$3$ & $2$ & $\ge 8$ & $\approx 47.59$ & Not found \\
$2$ & $3$ & $\ge 8$ & $\approx 83.02$ & Not found \\
$1$ & $4$ & $\ge 8$ & $\approx 201.57$ & Not found \\
\bottomrule
\end{tabular}
\end{table}

\begin{remark}
Among all weight-$5$ partitions, $S_{20}$ is the
\emph{only new sequence} (beyond the known Ap\'ery-type $(2,2)$
case) for which a finite recurrence has been established.
This suggests that the $(4,1)$ partition occupies a special
position in the hierarchy.
The remaining partitions may satisfy recurrences of order~$\ge 8$,
but our search with $30$ exact terms and polynomial degree
bounds up to $9$ did not produce them.
\end{remark}


%% ===================================================================
\section{Mirror Map Integrality}\label{sec:mirror}
%% ===================================================================

Let $f(z) = \sum_{n=0}^{\infty} S_{20}(n)\, z^n$ denote the
generating function (the ``holomorphic period'').
Following the Calabi--Yau paradigm \cite{LianYau1996,AESZ},
define the logarithmic period
\[
g(z) = \sum_{n=0}^{\infty} S_{20}(n)\, H(n)\, z^n
\]
where $H(n)$ encodes an appropriate harmonic-number combination,
and the \emph{mirror map}
\begin{equation}\label{eq:mirror}
q(z) = z\,\exp\!\bigl(g(z)/f(z)\bigr)
= z + \sum_{d=2}^{\infty} q_d\, z^d.
\end{equation}

\begin{theorem}[Mirror map integrality]\label{thm:mirror}
All coefficients $q_d$ of the mirror map~\eqref{eq:mirror}
associated with $S_{20}(n)$ satisfy $q_d \in \ZZ$ for
$d \le 16$.
\end{theorem}

\begin{proof}[Proof (computational)]
The coefficients $q_d$ for $d \le 16$ were computed using
exact rational arithmetic (no floating-point approximation)
and verified to be integers.
\end{proof}

\noindent
The computed values are:

\begin{table}[ht]
\centering
\caption{Mirror map coefficients $q_d$ for $d \le 16$.}
\label{tab:mirror}
\smallskip
\begin{tabular}{r r}
\toprule
$d$ & $q_d$ \\
\midrule
2 & 9 \\
3 & 165 \\
4 & 4\,110 \\
5 & 111\,075 \\
6 & 3\,316\,785 \\
7 & 104\,271\,733 \\
8 & 3\,421\,974\,692 \\
9 & 115\,918\,914\,756 \\
10 & 4\,027\,088\,171\,898 \\
11 & 142\,793\,489\,195\,634 \\
12 & 5\,149\,415\,166\,799\,466 \\
13 & 188\,353\,171\,046\,524\,999 \\
14 & 6\,973\,330\,284\,143\,733\,181 \\
15 & 260\,877\,511\,906\,858\,891\,334 \\
16 & 9\,848\,682\,801\,654\,949\,278\,015 \\
\bottomrule
\end{tabular}
\end{table}

\begin{conjecture}[Lian--Yau integrality]\label{conj:LY}
The mirror map coefficients $q_d \in \ZZ$ for \emph{all} $d \ge 1$.
\end{conjecture}

If true, this would place $S_{20}$ alongside the Ap\'ery numbers
and the $14$ Calabi--Yau operators catalogued by
Almkvist--van~Enckevort--van~Straten--Zudilin as sequences with
the Lian--Yau integrality property.


%% ===================================================================
\section{The Diagonal Representation Problem}\label{sec:diagonal}
%% ===================================================================

A key question in arithmetic combinatorics is whether a given
integer sequence is the \emph{diagonal} of a multivariate rational
function.  The Ap\'ery numbers are known to be the diagonal of
a four-variable rational function \cite{StrZud2004,ChristolDiag},
and this representation connects them to the period integrals of
a Calabi--Yau threefold.

For $S_{20}$, we pose:

\begin{conjecture}[Diagonal representation]\label{conj:diagonal}
There exists a rational function
$F \in \QQ(x_1, x_2, x_3, x_4, x_5)$ such that
\[
S_{20}(n) = [x_1^n x_2^n x_3^n x_4^n x_5^n]\,
\frac{1}{F(x_1, x_2, x_3, x_4, x_5)}.
\]
\end{conjecture}

\begin{remark}[Negative results]\label{rem:diag_negative}
We tested $17$ candidate rational functions of the form
$F = 1 - P(x_1,\ldots,x_5)$ where $P$ ranges over natural
generalisations of the known Ap\'ery diagonal.  In each case,
the diagonal of $1/F$ was computed through sufficiently many
terms and compared against $S_{20}(n)$; all~$17$ candidates failed.

In particular, the function
\[
F(x_1,\ldots,x_5) =
\frac{1}{1 - x_1(1-x_2)(1-x_3)(1-x_4)(1-x_5)
- x_1 x_2 x_3 x_4 x_5}
\]
was verified \emph{not} to have $S_{20}(n)$ as its diagonal:
an algebraic computation shows that its diagonal coefficients
equal $2^n$.
\end{remark}

\begin{proposition}[Existence, via Bostan--Lairez--Salvy \cite{BLS2015}]
\label{prop:diagonal_exists}
A diagonal representation of $S_{20}$ \emph{exists}.
\end{proposition}

\begin{proof}
By \cite[Theorem~3]{BLS2015}, a sequence $(a_n)$ is the
diagonal of a rational function if and only if it is D-finite
(satisfies a linear recurrence with polynomial coefficients).
Since $S_{20}$ satisfies the order-$5$ recurrence of
Theorem~\ref{thm:recurrence}, it is D-finite and hence a diagonal.

Moreover, the proof in \cite{BLS2015} is constructive:
using the Griffiths--Dwork reduction method, one can algorithmically
produce an explicit rational function whose diagonal is $S_{20}$.
However, the construction generically yields a function in
$9$ or more variables, and the challenge is to find the
\emph{minimal-variable} representation.
\end{proof}

\begin{remark}\label{rem:obstruction}
The $\tfrac{3}{4}$-well-poised structure identified in
Theorem~\ref{thm:hyp} may explain the difficulty of finding a
\emph{simple} diagonal.
The Ap\'ery diagonal \cite{Straub2014}
$A(n) = \diag\bigl[1/\bigl((1-x_1-x_2)(1-x_3-x_4) - x_1x_2x_3x_4\bigr)\bigr]$
arises from a fully well-poised ${}_4F_3$.
The partial well-poisedness of $S_{20}$ suggests that if a
\emph{simple} diagonal representation exists in $5$ variables,
it requires a qualitatively different rational function.
\end{remark}


%% ===================================================================
\section{Lean~4 Verification}\label{sec:lean}
%% ===================================================================

As a partial formal check, the identity~\eqref{eq:rec_n0}
(the recurrence at $n = 0$) has been compiled in the Lean~4
proof assistant:

\begin{verbatim}
def S20 (n : Nat) : Int :=
  ((Finset.range (n + 1)).sum
    (fun k => (Nat.choose n k)^4
            * (Nat.choose (n + k) k)))

theorem theorem20_inst0 :
    (-5412650858431135013634958175726842170573378411840)
      * S20 0
  + (-6600211789894833600749251782579095561783149274990400)
      * S20 1
  + (-29724234537629673550738669814459138431115401303206240)
      * S20 2
  + (-6675296886001563027617164081383167394996985596478240)
      * S20 3
  + (-272198721521932617277293245047721130052020296806560)
      * S20 4
  + (20478134952232355172884134183653971676016433020000)
      * S20 5
  = 0 := by decide
\end{verbatim}

\begin{remark}[Scope of the Lean verification]\label{rem:lean_scope}
The Lean~4 theorem \texttt{theorem20\_inst0} verifies a
\emph{single numerical identity}: the linear combination of
$S_{20}(0), \ldots, S_{20}(5)$ with the given integer
coefficients equals zero.
This confirms the recurrence at the point $n = 0$ but does
\emph{not} constitute a proof of the recurrence for all~$n$.
A full formal proof would require either formalising the
creative telescoping certificate or establishing a uniform
inductive argument, which remains an open problem.
\end{remark}


%% ===================================================================
\section{Open Questions}\label{sec:open}
%% ===================================================================

We collect the principal open problems arising from this work.

\begin{enumerate}
\item \textbf{Diagonal representation.}
Does there exist a rational function
$F \in \QQ(x_1,\ldots,x_5)$ whose diagonal is $S_{20}(n)$?
(Conjecture~\ref{conj:diagonal}.)
If so, does the associated hypersurface define a
Calabi--Yau fourfold?

\item \textbf{Lian--Yau integrality.}
Is $q_d \in \ZZ$ for all $d \ge 1$?
(Conjecture~\ref{conj:LY}.)
A proof would likely require a geometric construction
(e.g., a mirror pair of Calabi--Yau fourfolds).

\item \textbf{Full formal verification.}
Can the order-$5$ recurrence be proved for all $n$ in a
proof assistant, either by formalising the creative
telescoping certificate or by an independent method?

\item \textbf{Supercongruences.}
Do there exist $p$-adic supercongruences for $S_{20}$
analogous to those known for Ap\'ery numbers?
The $\tfrac{3}{4}$-well-poised structure may impose
constraints on which primes exhibit enhanced congruences.

\item \textbf{Modularity.}
Is the generating function $f(z) = \sum S_{20}(n) z^n$
a modular form or a period of a modular motive?
The order-$5$ recurrence operator should be analysed for
$G$-operator properties in the sense of Andr\'e.

\item \textbf{Remaining weight-$5$ partitions.}
What are the recurrence orders for the partitions
$(3,1)$, $(5,0)$, $(3,2)$, $(2,3)$, $(1,4)$?
Are they truly of order $\ge 8$, or do they require
higher-degree polynomial coefficients?

\item \textbf{Picard--Fuchs operator.}
The Picard--Fuchs ODE has singularities at
$t \approx 0.0232, -0.0692, -0.0736 \pm 0.0355i, -238.8$.
The last singularity is far from the origin.
Is it an apparent singularity, and does the operator
factorise over $\QQ(t)$?
\end{enumerate}


%% ===================================================================
%% References
%% ===================================================================

\begin{thebibliography}{99}

\bibitem{ABG2013}
G.~Almkvist, M.~Bogner, and J.~Guillera,
\emph{About a class of Calabi--Yau differential equations},
preprint, \texttt{arXiv:1310.6658}, 2013.

\bibitem{AESZ}
G.~Almkvist, C.~van~Enckevort, D.~van~Straten, and W.~Zudilin,
\emph{Tables of Calabi--Yau equations},
preprint, \texttt{arXiv:math/0507430}, 2005; updated 2010.

\bibitem{Apery1979}
R.~Ap\'ery,
\emph{Irrationalit\'e de $\zeta 2$ et $\zeta 3$},
Ast\'erisque \textbf{61} (1979), 11--13.

\bibitem{BLS2015}
A.~Bostan, P.~Lairez, and B.~Salvy,
\emph{Multiple binomial sums},
J.~Symbolic Comput.\ \textbf{80} (2017), 351--386;
\texttt{arXiv:1510.07487}.

\bibitem{ChristolDiag}
G.~Christol,
\emph{Diagonals of rational functions},
in: European Women in Mathematics (Luminy, 2013), CRM Proc.\ Lecture Notes.

\bibitem{LianYau1996}
B.~H.~Lian and S.-T.~Yau,
\emph{Mirror maps, modular relations and hypergeometric series~I},
preprint, \texttt{arXiv:hep-th/9507151}, 1995.

\bibitem{vanderPoorten1979}
A.~van~der~Poorten,
\emph{A proof that Euler missed \ldots\ Ap\'ery's proof of the
irrationality of $\zeta(3)$},
Math.\ Intelligencer \textbf{1} (1979), 195--203.

\bibitem{Straub2014}
A.~Straub,
\emph{Multivariate Ap\'ery numbers and supercongruences of
rational functions},
Algebra Number Theory \textbf{8} (2014), 1985--2008;
\texttt{arXiv:1401.0854}.

\bibitem{StrZud2004}
V.~Strehl,
\emph{Binomial identities --- combinatorial and algorithmic
aspects},
Discrete Math.\ \textbf{136} (1994), 309--346.

\bibitem{Zagier2009}
D.~Zagier,
\emph{Integral solutions of Ap\'ery-like recurrence equations},
in: Groups and Symmetries: From Neolithic Scots to John McKay,
CRM Proc.\ Lecture Notes \textbf{47}, AMS, 2009, pp.~349--366.

\end{thebibliography}

\end{document}
